% Generated draft: requires mathematical and source review.
\documentclass[11pt]{article}
\usepackage[margin=25mm]{geometry}
\usepackage{fontspec}
\setmainfont{Latin Modern Roman}
\usepackage{amsmath,amssymb,mathtools,unicode-math}
\setmathfont{Latin Modern Math}
\usepackage{enumitem,needspace,xurl,hyperref}
\setlength{\parindent}{0pt}
\setlength{\parskip}{5pt}
\setlength{\emergencystretch}{3em}
\newcommand{\N}{\mathbb{N}}
\newcommand{\Z}{\mathbb{Z}}
\newcommand{\Q}{\mathbb{Q}}
\newcommand{\R}{\mathbb{R}}
\newcommand{\C}{\mathbb{C}}
\newcommand{\F}{\mathbb{F}}
\newcommand{\A}{\mathbb{A}}
\newcommand{\PP}{\mathbb{P}}
\newcommand{\E}{\mathbb{E}}
\newcommand{\eps}{\varepsilon}
\newcommand{\dd}{\,\mathrm{d}}
\begin{document}
\section*{UnsolvedMath ID: 30004594}
\textbf{Source problem code:} catalogue entry
\subsection{Definitions and mathematical statement}
The supplied problem asks whether the survival threshold and the corresponding percolation threshold of the classical nearest-neighbor contact process satisfy
\[
\lambda_c<\lambda_p.
\]
The underlying graph, its dimension, and the meaning of the percolation threshold are not specified in the supplied record, so this comparison does not yet determine a unique mathematical statement.

For a conventional definition of the survival threshold, fix an infinite, locally finite graph \(G=(V,E)\) whose edges specify the nearest-neighbor relation, and a vertex \(o\in V\). With recovery rate normalized to \(1\) and infection rate \(\lambda\) per infected neighbor, the contact process \((\eta_t)_{t\geq0}\) takes values in \(\{0,1\}^{V}\). An infected vertex changes from \(1\) to \(0\) at rate \(1\), and an uninfected vertex \(x\) changes from \(0\) to \(1\) at rate
\[
\lambda\sum_{y:\{x,y\}\in E}\eta_t(y).
\]
Writing \(\eta_t^{\{o\}}\) for the infected set at time \(t\) when initially only \(o\) is infected, define the finite-seed global-survival threshold by
\[
\lambda_c(G,o)=\inf\left\{\lambda\geq0:
\mathbb P_{\lambda}^{G,\{o\}}\!\left(\eta_t^{\{o\}}\neq\varnothing\text{ for every }t\geq0\right)>0\right\},
\]
with \(\inf\varnothing=+\infty\). This is an explicitly stated conventional normalization and interpretation, whose applicability to the intended source formulation requires confirmation.

The requested comparison would then be \(\lambda_c(G,o)<\lambda_p\), with both thresholds referring to the same specified model and parameter convention. However, the record does not define the probability law or connectivity event underlying \(\lambda_p\). In particular, it does not say whether percolation concerns infected vertices in a stationary spatial configuration or another construction. It also does not specify the graph or the dimensions for which the inequality is requested. These missing choices prevent an unconditional, fully quantified formulation.

The supplied dated assessment identifies the classical nearest-neighbor case as open and says that the cited primary work resolves only a sufficiently spread-out model. That qualification is retained as supplied curation, not as an independently verified mathematical or literature claim.
\subsection{Short English statement}
Does the classical nearest-neighbor contact process have a survival threshold strictly below its corresponding percolation threshold? The graph, dimension, parameter convention, and definition of percolation must be confirmed before this becomes a fully specified conjecture. The supplied assessment labels the nearest-neighbor question open and distinguishes it from a sufficiently spread-out result.
\subsection{Sources}
\begin{itemize}
\item Local UnsolvedMath registry: permanent ID \(30004594\), source problem code \texttt{catalogue entry}, title \emph{Survival Versus Percolation Thresholds in the Contact Process}, in the Oberwolfach Reports Open Problems collection. The record supplies proposed year \(2021\) and names no proposer.
\item Supplied report citation: \emph{Spatial Networks and Percolation (2022)}, Oberwolfach Reports, pp.\ 149--224, DOI \texttt{10.4171/owr/2021/4}. \url{https://doi.org/10.4171/owr/2021/4}. The background describes these reports as documenting mathematical workshops and their open-problem discussions.
\item Supplied literature source: Rath and Valesin, spread-out contact process. \url{https://arxiv.org/abs/1912.09825}.
\item Supplied literature source: Liggett and Steif. \url{https://doi.org/10.1016/j.anihpb.2005.04.002}.
\item Literature assessment dated exactly \texttt{2026-08-22}: the supplied curation labels the classical nearest-neighbor strict inequality open and states that the cited primary work resolves only the sufficiently spread-out model. This is supplied dated curation, not independently verified; the record itself requests independent verification of time-sensitive status and cited claims.
\item The supplied statement-status label is \texttt{corrected\_verified}; its accompanying verification note concerns dollar-delimited formulas and preservation of supplied mathematical content. It does not establish mathematical correctness or resolve the missing hypotheses.
\end{itemize}
\end{document}
